% Finite multiple-root entry. Uses the existing joint-primitive bounds.
\subsection{Several root entries from one recovered contribution}
\label{sec:multiple-root-entry}

One recovered contribution can feed several available old roots and complete
root-information strings. The comparison below retains these uses jointly.
Let $S$ be uniform on 256 bits, initially independent of the public compression
oracle $h$ and auxiliary state, and retain $F$ and $\Gamma$ from
Section~\ref{sec:joint-continuation}. Each root request supplies a 32-byte
salt $r$ and a complete native 97-byte information string $I$. Its real value is
\[
 P_r=F(r,S),\qquad R_{r,I}=F(P_r,I\Vert\mathtt{01}).
\]
The parent contribution is used through these Extract calls; each resulting
PRK is used through its root-Expand interface. The controller retains the
original source ledger and any established local stops.

\paragraph{Descriptors and available data.}
The root-use descriptor is
$e(r,I)=(\Gamma(r),I)$, with exact byte equality. Equal descriptors retrieve
the same root entry across callers and resumptions. Indeed,
\eqref{eq:appendix-column-compression} shows that equal $\Gamma$ values give
equal Extract values for every $S$, and equal complete information strings
then give equal Expand values. Assign each newly encountered $\Gamma$ class
its next lifetime index $i$, preserving this index after a restart.
Numerical output coincidences between distinct classes keep their ordinary
shared-oracle distribution.

The controller computes salts, information strings and its selection from
its available state: executed preparation, independent coins, public-oracle
responses, unmarked roots, and the declared outputs of the marked schedule.
At most one fresh descriptor is marked before its first root evaluation.
Its root is installed in one retained schedule handle; every descriptor alias
retrieves that handle. All its uses belong to the same interface of
Section~\ref{sec:joint-continuation}, including its session salt, confirmation
and successor calls, complete traffic-information table and budgets. Other
root entries may be disclosed and used in later available computations,
including as salts. The retained handle supplies exactly that schedule's
operations. This is a rule about available computation; a locally known salt
can numerically coincide with a protected value.

Fix finite bounds $D,C,N_1,\ldots,N_C$ before execution: at most $D$ distinct
salt strings, at most $C\le D$ classes, and at most $N_i$ distinct complete
information strings in the $i$th class. These bounds hold in every comparison
experiment. All oracle-dependent preparation is executed and counted. Exact
repeats are cached. Write $q_S$ for the complete ordinary compression-query
budget in the column reduction, and $q_i$ for that in the $i$th row reduction.
These budgets include the whole controller, all other real primitives,
continuation and terminal test. An additional independent whole-hash oracle
may remain available throughout.

\begin{theorem}[Finite multiple-root entry]
\label{thm:multiple-root-entry}
Replace the real root interface by a private table
$\mathcal W(\Gamma(r),I)$ of independent uniform 32-byte entries, preserving
all descriptor aliases and the controller above. Every efficient available
output-bit test changes by at most
\begin{equation}
 \Delta_{\rm mr}=\min\!\left\{1,
 \varepsilon_{\rm Hjoint}(C,0,q_S,32)
 +\sum_{i=1}^{C}\varepsilon_{\rm Hrow}(N_i,q_i,98)\right\}.
 \label{eq:multiple-root-error}
\end{equation}
In the table experiment, an online marked fresh descriptor supplies the
independent uniform root required by
Theorem~\ref{thm:joint-message-leaf}, together with the complete interactive
controller. The marked descriptor may belong to a previously used class.
\end{theorem}
\begin{proof}
Compute $\Gamma(r)$ from available salt bytes and retain its class index.
Use one column call for the first salt in each new class; reuse its reply for
later equal classes. The descriptor identity makes this an exact rendering
of the real Extract interface. Lemma~\ref{lem:appendix-joint-hmac} replaces
these calls by $P_i=V(\Gamma(r))$ at the first cost in
\eqref{eq:multiple-root-error}. The remaining root computation consists of
HMAC rows under the independent private keys $P_i$.

Replace these rows in first-class-allocation order. At step $i$, presample
the hidden uniform row-challenge key and assign it to the $i$th new class
when that class is first encountered. Before that event the key has no use.
Classes already replaced use their private tables; later classes have locally
sampled independent keys and are evaluated with the same $h$. The controller
accesses the selected PRK through its row alone. Thus the real row branch
matches the preceding hybrid and the ideal row branch matches the next one.
Lemma~\ref{lem:appendix-hmac-instances} bounds this step by
$\varepsilon_{\rm Hrow}(N_i,q_i,98)$. If the class is never allocated, the
hybrids agree. Summing over the fixed class budget establishes the bound.

Distinct classes remain distinct reduction slots even when their sampled
PRK bytes coincide. In a real row branch, local evaluation under another
class's key uses the same $h$ and automatically preserves any resulting
coincidence; these ordinary queries are included in $q_i$. Thus the argument
uses the full independent-key distribution.

For each ideal row, $I\mapsto I\Vert\mathtt{01}$ is injective. Its root
interface is therefore exactly the table $\mathcal W(\Gamma(r),I)$.
At a fresh descriptor, its value is independent of all previously used table
entries and the oracle. Presample the designated uniform root, withhold it
during the executed prefix, and assign it when the controller marks that
fresh descriptor. All other entries retain their independent sampling and
all aliases retrieve the retained handle. The selection uses the available
prefix, so this deferred assignment has exactly the table distribution.
It embeds the whole controller in the existing root experiment. The row
hybrids already covered every class, so target selection adds no class guess.
\end{proof}

\paragraph{Explicit query accounting.}
Let $B$ bound all ordinary compression work outside this entry service,
including its entire downstream schedule and test. Cache the two public
$\Gamma$ computations per distinct salt. With
$N=\sum_{k=1}^{C}N_k$, valid conservative choices are
\begin{equation}
 q_S=B+2D+5N,\qquad
 q_i=B+2D+5\sum_{k=i+1}^{C}N_k.
 \label{eq:multiple-root-budgets}
\end{equation}
A 98-byte HMAC evaluation uses five compression calls, including its two key
pads. In the column comparison every root row is real; at row step $i$, the
later rows are real and the earlier rows are tables. External class scheduling
accounts for $2D$, including new salts that share an existing class. The
column primitive's own public pad work remains included in its established
bound. $B$ is a uniform bound over these complete derived experiments.

\begin{corollary}[First message after a marked root entry]
\label{cor:multiple-root-message}
For the available controller above and its marked schedule, use the
first-message interface and available qualification predicate of Corollary~\ref{cor:joint-first-message}. If $Z$
is its qualified signed message score, with zero on runs having no marked
challenge, then
\begin{equation}
 |\mathbb E Z|\le
 \min\{1,\,2\Delta_{\rm mr}+2\Delta_{\rm sched}
                      +2\varepsilon_{\rm AEAD}\}.
 \label{eq:multiple-root-message}
\end{equation}
The schedule and AEAD reductions include the complete entry controller and
its shared-oracle work. If the contribution itself comes from the protected
ContextBound table of Lemma~\ref{lem:cb-hidden-component}, its additional
signed-score cost is $2\Phi_{256}(d,Q_H)$. For a candidate whose stopped execution realizes this entry controller,
it replaces the single-entry $2\Delta_{\rm init}$ term by
$2\Delta_{\rm mr}$ in the corresponding continuation bound; the original accounting terms keep their
established definitions.
\end{corollary}
\begin{proof}
Fix each application challenge bit for the entry comparison and average
afterwards, as in the retained continuation. The randomized output test
with success probability $(Z+1)/2$ converts
Theorem~\ref{thm:multiple-root-entry}'s bit gap into a score gap of at most
$2\Delta_{\rm mr}$. Its ideal marked root initializes the entire controller
of Corollary~\ref{cor:joint-first-message}, which supplies the remaining two
terms. Apply the same conversion to the ContextBound comparison when present.
\end{proof}
